\chapter*{Introduction}
\phantomsection
\addcontentsline{toc}{chapter}{Introduction}

\section{Contexte général}
Au cours des dernières années, le domaine du sport a connu une transformation significative sous l’impulsion des avancées technologiques, notamment dans les domaines de l’intelligence artificielle et de la vision par ordinateur. Ces technologies ont progressivement redéfini la manière dont les performances sportives sont analysées, interprétées et exploitées, tant dans un cadre professionnel qu’amateur.

Traditionnellement, l’analyse des matchs et la gestion des événements de jeu reposaient en grande partie sur l’intervention humaine. Qu’il s’agisse de l’arbitrage, du suivi du score ou de l’analyse des performances, ces tâches impliquaient une observation continue et une prise de décision en temps réel, souvent sujette à des limitations liées à la fatigue, à la subjectivité ou à la complexité des situations de jeu. Dans ce contexte, l’automatisation de ces processus apparaît comme un enjeu majeur pour améliorer la fiabilité, la précision et la continuité de l’expérience sportive.

Parallèlement, l’essor des systèmes embarqués, des caméras à haute résolution et des capacités de calcul a ouvert la voie à de nouvelles approches basées sur l’analyse vidéo. Ces approches permettent de capturer, traiter et interpréter des flux visuels en temps réel, offrant ainsi la possibilité de détecter automatiquement des objets, de suivre leurs mouvements et d’identifier des événements significatifs au sein d’un match.

Dans cette dynamique, les solutions dites de \textit{« sport-tech »} émergent comme un levier d’innovation, visant à digitaliser et optimiser les processus liés à la gestion et à l’analyse des activités sportives. Elles s’inscrivent dans une volonté plus large de rendre les données sportives plus accessibles, exploitables et pertinentes pour différents acteurs, qu’il s’agisse des joueurs, des entraîneurs ou des organisateurs.

C’est dans ce cadre général que s’inscrit le présent travail, qui explore l’apport des techniques d’intelligence artificielle et de vision par ordinateur dans l’automatisation de l’analyse des matchs, avec un intérêt particulier pour la détection des événements de jeu et la génération automatique du score.


\section*{Problématique}

Malgré les avancées technologiques observées dans le domaine de l’analyse sportive, la gestion du score et la détection des événements de jeu reposent encore majoritairement sur des interventions humaines ou sur des systèmes semi-automatisés. Cette dépendance introduit un certain nombre de limites, notamment en termes de fiabilité, de continuité et de précision, en particulier dans des contextes de jeu rapides et dynamiques.

Dans des disciplines telles que le tennis ou le padel, où les échanges sont courts mais intenses et où les décisions doivent être prises en temps réel, l’identification des événements clés --- tels que les points gagnants, les fautes ou les sorties de balle --- constitue un défi complexe. Ce défi est accentué par plusieurs facteurs critiques :

\begin{itemize}
    \item La vitesse élevée de la balle et les variations brusques de trajectoire ;
    \item Les phénomènes d’occlusion (joueurs masquant la balle ou le terrain) ;
    \item Les conditions environnementales variables (éclairage, angles de vue, qualité du flux vidéo).
\end{itemize}

Par ailleurs, les solutions existantes présentent souvent des limites en termes d’accessibilité, de coût ou de complexité d’installation, ce qui restreint leur adoption à grande échelle. Les systèmes les plus avancés nécessitent généralement des infrastructures lourdes, tandis que les approches plus légères peinent à garantir un niveau de précision suffisant.

Dans ce contexte, la problématique centrale de ce travail consiste à concevoir un système capable d’automatiser, de manière fiable et en temps réel, la détection des événements de jeu à partir de flux vidéo, en vue de produire un score précis sans intervention humaine. Il s’agit de répondre aux contraintes de robustesse face aux incertitudes visuelles, de maintenir une performance temporelle compatible avec le temps réel, et d’assurer une cohérence logique dans l’interprétation des données.

\subsection*{Question de recherche}
Ainsi, la question à laquelle ce projet cherche à répondre peut être formulée comme suit :
\begin{quote}
    \textit{« Comment concevoir un système intelligent, basé sur l’analyse vidéo, capable de détecter et d’interpréter automatiquement les événements de jeu afin de générer un score fiable, en temps réel et dans des conditions réelles d’utilisation ? »}
\end{quote}

\section{Solution proposée}

Afin de répondre à la problématique identifiée, ce travail propose la conception et la mise en œuvre d’un système intelligent de scoring automatique basé sur l’analyse de flux vidéo. L’objectif principal est de détecter, interpréter et exploiter les événements de jeu en temps réel, afin de générer un score fiable sans intervention humaine.

La solution adoptée repose sur une architecture modulaire, structurée en plusieurs niveaux de traitement complémentaires. Dans un premier temps, un pipeline de vision par ordinateur est mis en place pour assurer la détection et le suivi des éléments principaux du jeu, notamment les joueurs et la balle. Cette étape s’appuie sur des modèles d’apprentissage profond permettant d’extraire des informations pertinentes à partir des images vidéo.

Ensuite, les données issues de la détection sont enrichies par une phase de modélisation spatiale du terrain. Cette étape permet de contextualiser les positions et les trajectoires observées, en les projetant dans un référentiel cohérent, facilitant ainsi l’interprétation des interactions entre les éléments du jeu.

Sur cette base, un module de raisonnement est introduit afin d’identifier les événements clés du match, tels que les échanges, les fautes ou les points marqués. Ce module s’appuie sur une logique formelle, implémentée à travers une machine à états finis, permettant de modéliser le déroulement du jeu et d’assurer la cohérence des décisions prises.

Enfin, les informations extraites sont utilisées pour générer automatiquement le score en temps réel, ainsi que pour produire des statistiques et des visualisations exploitables. L’ensemble du système est conçu pour fonctionner dans des conditions proches du réel, en tenant compte des contraintes de performance, de robustesse et de latence.

Ainsi, la solution proposée se distingue par l’intégration cohérente de techniques de vision par ordinateur, de modélisation géométrique et de raisonnement logique, au sein d’un pipeline unifié dédié à l’automatisation de l’analyse des matchs sportifs.

\section*{Organisation du rapport}

Le présent rapport est structuré en trois chapitres principaux, organisés de manière à présenter progressivement le contexte du projet, la solution proposée, ainsi que sa mise en œuvre et son évaluation.

Le premier chapitre, intitulé « Positionnement du projet et état de l’art », introduit le cadre général du travail. Il présente le contexte dans lequel s’inscrit le projet, précise la problématique étudiée et décrit les principales solutions existantes dans le domaine. Une analyse comparative est ensuite réalisée afin de mettre en évidence les limites des approches actuelles et de justifier les choix technologiques adoptés.

Le deuxième chapitre, « Architecture de la solution », est consacré à la conception du système proposé. Il détaille l’architecture globale ainsi que les différents modules qui la composent. Ce chapitre décrit notamment le pipeline de traitement vidéo, les méthodes de détection et de suivi des joueurs et de la balle, la modélisation du terrain, ainsi que les mécanismes de projection géométrique. Il présente également le moteur de scoring automatique, les techniques de génération des statistiques, ainsi que les aspects liés aux performances et à l’optimisation du système.

Le troisième chapitre, « Implémentation, tests et évaluation », porte sur la réalisation pratique de la solution. Il expose l’environnement de développement, les outils utilisés et l’organisation du code. Les différentes étapes de développement sont décrites, suivies d’une présentation de la méthodologie de test et des scénarios de validation. Enfin, ce chapitre propose une analyse des résultats obtenus à travers des métriques quantitatives, tout en discutant les points forts, les limites du système et les perspectives d’amélioration.

Le rapport se conclut par une conclusion générale qui synthétise les contributions du projet et ouvre sur des pistes de recherche et d’évolution futures.